\documentclass[12pt, a4paper]{article}

\usepackage[a4paper, top=2cm, bottom=2cm, left=2cm, right=2cm]{geometry}

\usepackage{tgheros}
\renewcommand{\familydefault}{\sfdefault}
\usepackage[T1]{fontenc}
\usepackage{xcolor}

\usepackage{fancyhdr}
\pagestyle{fancy}

\usepackage{tocloft}
\usepackage{hyperref}

\usepackage{graphicx}
\usepackage{float}

\usepackage{booktabs}

\hypersetup{
	colorlinks=true,
	linkcolor=black,
	urlcolor=blue
}

\definecolor{CorFundo}{HTML}{f23847}

\renewcommand{\contentsname}{\textbf{\LARGE \fontsize{34}{46}\selectfont Tabela de conteúdos}}

\fancypagestyle{sumariostyle}{
	\fancyhf{}
	\renewcommand{\headrulewidth}{0pt}
	\renewcommand{\footrulewidth}{1.3mm}
	\renewcommand{\footrule}{{\color{CorFundo}\hrule width\headwidth height\footrulewidth}}
	\fancyfoot[L]{\textcolor{CorFundo}{Revisão 1.0 --- Outubro de 2026}}
}

\begin{document}
	\begin{titlepage}
		\color{white}
		\pagecolor{CorFundo}
		
		\begin{flushleft}
			\noindent \small \textbf{DFC-8} Um Microcontrolador por IFBaiano -- Campus Guanambi \\
			\rule{\linewidth}{1.3mm}
			\vspace{1.5cm}
			
			{\huge \textbf{DFC-8 Datasheet} \\}
			\vspace{0.2cm}
			{\LARGE Um Microcontrolador \\}
			\vspace{0.2cm}
			{\LARGE por IFBaiano -- Campus \\}
			\vspace{0.2cm}
			{\LARGE Guanambi}
			
		 	\vfill
		 	
			{\color{white}\rule{\linewidth}{1.3mm}} \\[0.2cm]
			\textcolor{white}{\small Revisão 1.0 --- Outubro de 2026}
					 	
		\end{flushleft}
	\end{titlepage}
	
	\pagenumbering{roman}
	
	\clearpage
	
	\pagecolor{white}
	\color{black}
	
	\pagestyle{sumariostyle}
	\addtocontents{toc}{\protect\thispagestyle{sumariostyle}}
	\tableofcontents
	\newpage
	\thispagestyle{sumariostyle}
	
	\clearpage

	\pagenumbering{arabic}
	\pagestyle{fancy}
	
	\fancyhf{}
	\renewcommand{\headrulewidth}{0.4pt}
	\renewcommand{\footrulewidth}{0.4pt}
	
	\section{Introdução}
	Os microcontroladores conectam o mundo do software ao mundo do hardware. Eles permitem que os desenvolvedores escrevam softwares que interagem com o mundo físico da mesma maneira determinística e precisa a nível de ciclo como a lógica digital. Eles ocupam o canto inferior esquerdo do espaço de preço/desempenho, vendendo dez vezes mais que seus irmãos mais potentes. Eles são os pilares que impulsionam a transformação digital do nosso mundo.
	
	\subsection{Arquitetura}
	O DFC-8 é um microcontrolador com arquitetura baseada em 8051. Ele é focado no controle do fluxo de dados de pilhas e filas, permitindo o controle total da passagem de dados pelas estruturas de forma dinâmica. Possuindo dois registradores principais, onde um é responsável pelo controle e outro pelo status das estruturas.
	
		\begin{figure}[H]
			\centering
			\includegraphics[width=0.7\linewidth]{8051-arch}
			\caption{MCS-51 Block Diagram}
			\label{fig:8051-arch}
		\end{figure}
	
	A arquitetura do DFC-8 foi otimizada para minimizar o overhead de processamento na manipulação de dados sequenciais. Através da gestão dinâmica dos ponteiros de memória e da monitorização contínua dos limites das estruturas, o microcontrolador executa operações de inserção e remoção com eficiência de ciclo único de instrução.
	
	\section{Mapa de Endereços}
	
	O mapa de endereços a seguir detalha o mapeamento de memória para os registradores dedicados ao gerenciamento de estruturas de dados do DFC-8. O acesso a esses registradores é feito através de um barramento de endereçamento de 8 bits.
	
	\begin{table}[H]
		\centering
		\caption{Mapa de Endereços dos Registradores de Dados}
		\vspace{0.2cm}
		\renewcommand{\arraystretch}{1.3}
		\begin{tabular}{@{}lllc@{}}
			\toprule
			\textbf{Endereço} & \textbf{Registrador} & \textbf{Descrição} & \textbf{Acesso} \\ \midrule
			\texttt{0x40}     & \texttt{DATA0\_CTRL}   & Registrador de controle de pilhas e filas & R/W    \\ 
			\texttt{0x60}     & \texttt{DATA0\_STATUS} & Registrador de status (estado) das estruturas & RO     \\ \bottomrule
		\end{tabular}
		\label{tab:mapa_enderecos}
	\end{table}
	
	\vspace{0.3cm}
	\noindent
	\textbf{Tipos de Acesso:}
	\begin{itemize}
		\item \textbf{R/W} (\textit{Read/Write}): Leitura e Escrita permitidas.
		\item \textbf{RO} (\textit{Read-Only}): Apenas leitura. Modificações são gerenciadas exclusivamente pelo hardware.
	\end{itemize}
	
	\vspace{0.5cm}
	
\subsection{DATA0\_CTRL}

O registrador \texttt{DATA0\_CTRL} controla o fluxo de dados da pilha e da fila.

	\begin{table}[H]
		\centering
		\caption{Descrição dos bits do registrador DATA0\_CTRL}
		\vspace{0.2cm}
		\renewcommand{\arraystretch}{1.5}
		\setlength{\arrayrulewidth}{1pt} % Deixa as linhas da tabela mais grossas/escuras
		\begin{tabular}{|c|p{10cm}|c|}
			\hline
			\textbf{Bits} & \textbf{Descrição} & \textbf{Tipo} \\ \hline
			7   & \textbf{FIFO\_BLOCK:} Bloqueia o fluxo de dados da fila. & RW \\ \hline
			6   & \textbf{LIFO\_BLOCK:} Bloqueia o fluxo de dados da pilha. & RW \\ \hline
			5:2 & Reservado. & -  \\ \hline
			1   & \textbf{FIFO\_UNLOCK:} Libera o fluxo de dados da fila. & RW \\ \hline
			0   & \textbf{LIFO\_UNLOCK:} Libera o fluxo de dados da pilha. & RW \\ \hline
		\end{tabular}
		\label{tab:data0_ctrl}
	\end{table}

\vspace{0.5cm}

\subsection{DATA0\_STATUS}

Esse é um registrador de leitura, e indica se a fila/pilha está cheia ou não.

	\begin{table}[H]
		\centering
		\caption{Descrição dos bits do registrador DATA0\_STATUS}
		\vspace{0.2cm}
		\renewcommand{\arraystretch}{1.5}
		\setlength{\arrayrulewidth}{1pt}
		\begin{tabular}{|c|p{10cm}|c|}
			\hline
			\textbf{Bits} & \textbf{Descrição} & \textbf{Tipo} \\ \hline
			7   & \textbf{FIFO\_FULL:} Quando a fila está cheia - 1. Quando ela ainda não encheu - 0. & RO \\ \hline
			6   & \textbf{LIFO\_FULL:} Quando a pilha está cheia - 1. Quando ela ainda não encheu - 0. & RO \\ \hline
			5:0 & Reservado. & -  \\ \hline
		\end{tabular}
		\label{tab:data0_status}
	\end{table}
	
	\clearpage
	\thispagestyle{empty}
	\pagecolor{CorFundo}
	\color{white}
	
	\mbox{}
		
		
	{\color{white}\rule{\linewidth}{1.3mm}} \\[0.2cm]
		
	\vspace*{\fill}
	
	\begin{flushleft}
		\begin{flushleft}
			{\LARGE \textbf{IFBaiano -- Campus Guanambi}} \\
		\end{flushleft}
		{\color{white}\rule{\linewidth}{1.3mm}} \\[0.2cm]
		\small Revisão 1.0 --- Outubro de 2026
	\end{flushleft}
\end{document}